% =====================================================================
%  simulation_methodology.tex
%  Replacement for the "Computer Simulation Methodology" section.
%  Updated for RadCluster_2_0: the SIA master equation is split into two
%  coupled loop-character populations — glissile c^{(111)}_n and sessile
%  c^{(100)}_n — consistent with the production C++ solver (rhs_case2 /
%  rhs_bin_moment) and the ½⟨111⟩→⟨100⟩ conversion edges (P9).
%
%  Requires: amsmath, amssymb, booktabs, natbib (or compatible \citep/\cite).
%  Macros assumed from the parent paper: \popset, \Vset, \Eset, \RAG.
%  \input this file in place of the old section.
% =====================================================================

\section{Computer Simulation Methodology}
\label{sec:simulations}

For \textsc{EUROFER-97} under neutron irradiation, the only resolved
trapped solute is helium:
\begin{equation}
   N_g = 1, \qquad g_1 = \mathrm{He}, \qquad
   \boldsymbol{\ell} = (\ell),
   \label{eq:eurofer_gas_list}
\end{equation}
so the gas reservoir $c^g_1(t)\equiv c_h(t)$ and the composition
vector reduce to the scalar helium occupancy $\ell\in\mathbb{Z}_{\geq 0}$.

\paragraph{Active populations}
\textsc{EUROFER-97} carries two self-interstitial (SIA) loop populations of
distinct Burgers-vector character on the positive polarity, together with a
single vacancy population on the negative polarity:
\begin{align}
   \popset_{-}^{\textsc{Eur}} &= \{\,\mathrm{bulk}\,\},
   \quad N_P^V = 1,
   \label{eq:eurofer_vac_pop}\\
   \popset_{+}^{\textsc{Eur}} &= \{\,\mathrm{bulk\text{-}111},\;\mathrm{bulk\text{-}100}\,\},
   \quad N_P^I = 2.
   \label{eq:eurofer_sia_pop}
\end{align}
The positive polarity resolves two distinct SIA loop characters as separate
populations: the glissile $\tfrac{1}{2}\langle 111\rangle$ loops
($\mathrm{bulk\text{-}111}$), which are 3D-mobile clusters for $n<4$ and
1D-glissile prismatic loops for $n\geq 4$, and the sessile $\langle 100\rangle$
loops ($\mathrm{bulk\text{-}100}$), which form during microstructural evolution
and act as essentially immobile sinks. Within each loop population the 3D/1D
mobility distinction is encoded in size-dependent rate kernels; the two
characters, by contrast, are tracked as separate population labels and are
coupled by the $\tfrac{1}{2}\langle 111\rangle\!\to\!\langle 100\rangle$
conversion edges introduced below. Grain boundaries, the dislocation network,
and MX/M$_{23}$C$_6$ precipitates are treated as fixed unresolved sinks.
On the negative polarity a single vacancy population is retained and its label
is suppressed, while on the positive polarity the loop-character label is kept
explicit:
\begin{equation}
   c^I_{n,\alpha}(t) \;\rightarrow\; c^{(111)}_n(t),\,c^{(100)}_n(t),
   \quad
   c^V_{m,(\ell),\beta}(t) \;\rightarrow\; c_{m,\ell}(t),
   \quad
   c^g_1(t) \;\rightarrow\; c_h(t).
   \label{eq:eurofer_variable_reduction}
\end{equation}

\paragraph{Admissible vertex set}
\begin{align}
   \Vset_{-}^{\textsc{Eur}}
   &= \bigl\{\,(-1,m,\mathrm{bulk},(\ell))\,:\,
              1\leq m\leq N_V,\;0\leq\ell\leq\ell_{\max}(m)\,\bigr\},
   \label{eq:eurofer_vset_minus}\\[2pt]
   \Vset_{+}^{\textsc{Eur}}
   &= \bigl\{\,(+1,n,p,\mathbf{0})\,:\,
              p\in\{\mathrm{bulk\text{-}111},\mathrm{bulk\text{-}100}\},\;1\leq n\leq N_I\,\bigr\},
   \label{eq:eurofer_vset_plus}\\[2pt]
   \Vset_{0}^{\textsc{Eur}}
   &= \{\,g_1=\mathrm{He}\,\},
   \label{eq:eurofer_vset_zero}
\end{align}
comprising the size-resolved glissile $\tfrac{1}{2}\langle 111\rangle$ and
sessile $\langle 100\rangle$ SIA loop populations, size-resolved voids
($\ell=0$) and helium bubbles ($\ell\geq 1$), and the single matrix He
reservoir. The sessile $\langle 100\rangle$ vertices are admissible only for
$n\geq n_{\mathrm{conv}}$ (the conversion onset size, Eq.~\eqref{eq:driving-force}).
The maximum helium content obeys the saturation rule
\begin{equation}
   \ell_{\max}(m) \;=\; \bigl\lfloor \alpha_{\mathrm{He}}\,m^{2/3}\bigr\rfloor,
   \qquad \alpha_{\mathrm{He}} \in [1.5,2.0].
   \label{eq:ellmax_eurofer}
\end{equation}

\paragraph{Primary populations}
Equations~\eqref{eq:eurofer_vset_minus}--\eqref{eq:eurofer_vset_zero}
yield three defect classes evolved by the master equations, with the SIA class
resolved into two coupled loop-character populations:
(1a)~\textbf{$\tfrac{1}{2}\langle 111\rangle$ SIA loops} $c^{(111)}_n(t)$,
$n=1,\ldots,N_I$: glissile prismatic loops for $n\geq 4$ and 3D-mobile clusters
for $n=1,2,3$, fed by the cascade source and carrying no helium;
(1b)~\textbf{$\langle 100\rangle$ SIA loops} $c^{(100)}_n(t)$,
$n\geq n_{\mathrm{conv}}$: sessile prismatic loops generated \emph{only} by
$\tfrac{1}{2}\langle 111\rangle\!\to\!\langle 100\rangle$ conversion (no cascade
source, no self-coalescence), carrying no helium;
(2)~\textbf{vacancy clusters / cavities} $c_{m,\ell}(t)$,
$m=1,\ldots,N_V$, $\ell=0,\ldots,\ell_{\max}(m)$: spherical cavities, pure
voids when $\ell=0$ and helium bubbles when $\ell\geq 1$;
(3)~\textbf{free helium} $c_h(t)$, the single gas-reservoir variable.
The signed lattice-defect contents of both SIA characters are positive,
\begin{equation}
   q\bigl(I^{(111)}_n\bigr) = +n,
   \qquad
   q\bigl(I^{(100)}_n\bigr) = +n,
   \qquad
   q(V_m\mathrm{He}_\ell) = -m,
   \qquad
   q(\mathrm{He}) = 0,
   \label{eq:q_eurofer}
\end{equation}
as required by the conservation law of Section~\ref{sec:FP_balance}; the
conversion edges relabel SIA content between the two positive populations and
therefore leave the total interstitial inventory unchanged.

All pairwise reaction rate constants are obtained from the Smoluchowski
capture model in atomic-fraction units. For a mobile point defect $\alpha$
absorbed by an immobile spherical cluster of $m$ defects the diffusion-limited
form reduces to
\begin{equation}
   \mathcal{K}^{\mathrm{3D}}_{\alpha,m}
   = 8\pi\,\xi(m)\,\omega_\alpha \;[\mathrm{s}^{-1}],
\end{equation}
with the void, dislocation-loop, recombination, and mixed 1D/3D-glide capture
constants ($\mathcal{K}^{\mathrm{void}}$, $\mathcal{K}^{\mathrm{loop}}_{i,n}$,
$\mathcal{K}_{iv}$, $\mathcal{K}^{\mathrm{mixed}}_{n,m}$), the 1D-glide reaction
frequencies $\omega^{\mathrm{1D}}_n$, and the correlated 1D versus 3D
sink-strength limits derived in detail in Section~\ref{sec:reactions}. Thermal
emission of species $\alpha$ from a size-$n$ cluster is fixed by detailed
balance against the capture rate,
\begin{equation}
   \mathcal{E}_\alpha(n)
   = \mathcal{K}_{\alpha}(n-1)\,
     \frac{1}{\Omega}\exp\!\left(-\frac{E_b^\alpha(n)}{k_BT}\right),
\end{equation}
with the corresponding loop-emission, vacancy/helium bubble-emission, and
trap-mutation expressions also collected in Section~\ref{sec:reactions}, where
Table~\ref{tab:rate_params} lists the rate-constant prefactors, sink
parameters, and state-space controls used throughout.

\paragraph{Process classes P1--P9}
The admissible edges of $\RAG^{\textsc{Eur}}$
($\chi_{ab\rightarrow c}=1,\ \&\ \chi_{a\rightarrow b}=1$) partition into nine
process classes, each mapped to an abstract edge class in
Section~\ref{sec:reaction_class}. Because the SIA class comprises two
loop-character populations, the inter-population transfer edge (P9) is active
and carries the $\tfrac{1}{2}\langle 111\rangle\!\to\!\langle 100\rangle$
conversion (Dudarev unary reorientation, Marian junction nucleation, and Marian
absorption growth); the defect production edge is treated as a source term that
feeds the $\tfrac12\langle111\rangle$ population alone.
\begin{center}
\renewcommand{\arraystretch}{1.1}
\begin{tabular}{@{}lll@{}}
\toprule
\textbf{Process} & \textbf{Abstract edge class} & \textbf{Rate kernel} \\
\midrule
P1: V--SIA recombination & cross-polarity (recombination) & $\mathcal{K}_{iv}$ \\
P2: cavity absorption of mono-$\alpha$ & growth / shrinkage & $\mathcal{K}^{\mathrm{cav}}_{\alpha,m}$ \\
P3: loop absorption of mono-$\alpha$ & growth / shrinkage & $\mathcal{K}^{\mathrm{loop}}_{\alpha,n}$ \\
P4: absorption at fixed sink & sink edge & $\mathcal{D}_\alpha$ \\
P5: thermal monomer emission & dissociation & $\varepsilon_\alpha$ \\
P6: 1D-gliding SIA cluster $\to$ cavity & cross-polarity & $\mathcal{K}^{\mathrm{eff}}_{n,m}$ \\
P7: trap mutation & gas-pressure-driven & $\Gamma_{\mathrm{TM}}(m,\ell)$ \\
P8: radiation re-solution & solute detrapping & $\Gamma_{\mathrm{res}}(m,\ell)$ \\
P9: $\tfrac{1}{2}\langle 111\rangle\!\to\!\langle 100\rangle$ conversion & inter-population transfer & $\Gamma_{\mathrm{uni}},\,\varphi_{\rm junc},\,K^{\rm abs}$ \\
\bottomrule
\end{tabular}
\end{center}

\paragraph*{Mobility cutoffs}
Mobile clusters are bounded by size limits from Table~\ref{tab:rate_params}:
$\tfrac12\langle111\rangle$ SIA clusters are mobile (3D or 1D) for
$1\leq n\leq i_{\rm mobile}$; larger $\tfrac12\langle111\rangle$ clusters are
sessile dislocation loops acting only as fixed sinks. The $\langle100\rangle$
loops are sessile at every size ($D^{100}_n\equiv 0$) and grow or shrink only
through single-monomer exchange. Vacancy clusters are 3D mobile for
$1\leq m\leq v_{\rm mobile}$; larger cavities grow or shrink only through
single-monomer absorption or emission. Under the mono-defect simplification
($i_{\rm mobile}=v_{\rm mobile}=1$), the coalescence rate constants reduce to
\begin{align}
    \mathcal{K}_{n,1}^{ii}\big|_{n'=1}
    &= \mathcal{K}^{\mathrm{loop}}_{i,n}
     = A_{\mathrm{loop}}\,n^{1/2}\,Z_i^{\mathrm{loop}}\,\omega_i^{\mathrm{eff}},
    \label{eq:monodef_i}\\[4pt]
    \mathcal{K}_{m,1}^{vv}\big|_{m'=1}
    &= \mathcal{K}^{\mathrm{cav}}_{v,m}
     = A_{\mathrm{sph}}\,m^{1/3}\,\omega_v^{\mathrm{eff}},
    \label{eq:monodef_v}\\[4pt]
    \mathcal{K}_{m=1,n=1}^{vi}
    &= \mathcal{K}_{iv}
     = 4\sqrt{3}\,\pi\,(\omega_i^{\mathrm{eff}}+\omega_v^{\mathrm{eff}}).
    \label{eq:monodef_iv}
\end{align}

%===================================================================
\paragraph{Reaction Processes (general coalescence formulation)}
\label{sec:reaction_class}
%===================================================================
The most complete description allows any two mobile clusters to react. The
intra- and cross-polarity channels are $\tfrac12\langle111\rangle$ SIA--SIA
coalescence ($I_n + I_{n'} \to I_{n+n'}$), vacancy--vacancy coalescence
($V_m + V_{m'} \to V_{m+m'}$), and V--I annihilation
($V_m + I_n \to V_{m-n},\,I_{n-m}$, or $\varnothing$). All three share the
diffusion-limited two-mobile-cluster form
\begin{equation}\label{eq:K_ii}
   \mathcal{K}_{a,b}
   = 8\pi\,(\xi_a + \xi_b)\,
     (\omega_a^{\mathrm{eff}} + \omega_b^{\mathrm{eff}}),
   \qquad \xi_a = (3a/8\pi)^{1/3},
\end{equation}
restricted to mobile sizes ($\leq i_{\rm mobile}$, $v_{\rm mobile}$); the
channel-specific constants and size restrictions are given in
Section~\ref{sec:reactions}. He--He coalescence is excluded because the
interstitial He--He binding energy in bcc Fe is negligible ($\lesssim 0.05$~eV).

The abstract master equation is assembled algorithmically by looping over every
admissible edge $r\in\Eset^{\textsc{Eur}}$, computing the flux $J_r$, and
accumulating the signed contribution $S_{ar}J_r$ into each affected population
$a$. Table~\ref{tab:gw_assembly} lists every edge class, its reaction, its flux,
and its non-zero stoichiometric contributions to the four EUROFER-97 population
equations $\dot{c}^{(111)}_n$, $\dot{c}^{(100)}_n$, $\dot{c}_{m,\ell}$, and
$\dot{c}_h$. Where a process connects adjacent nodes on a ladder (e.g.\ cavity
or sessile-loop growth/shrinkage), the gain term at $a$ and the loss term at $a$
appear on separate rows; the net effect is the tridiagonal ladder operator
visible in the boxed equations. Figure~\ref{fig:fm_network} shows the full
EUROFER-97 reaction network as a directed graph, with the nine process classes
color-coded on the arrows and the graph-walker formula annotated in the upper
center.

Define the total cavity concentration at vacancy size $m$ and the combined
fixed-sink rate:
\begin{equation}\label{eq:cbar}
    \bar{c}_m(t) \equiv \sum_{\ell=0}^{\ell_{\max}(m)} c_{m,\ell}(t),
    \qquad
    \mathcal{D}_\alpha
    = \mathcal{D}_\alpha^d + \mathcal{D}_\alpha^{gb}
    + \mathcal{D}_\alpha^p.
\end{equation}
The equations below are written in the general (coalescence) form; the
mono-defect simplification is recovered by setting
$i_{\rm mobile}=v_{\rm mobile}=1$ (Eqs.~\ref{eq:monodef_i}--\ref{eq:monodef_iv}).

%-------------------------------------------------------------------
\paragraph{Defect Class 1a: glissile $\tfrac12\langle111\rangle$ SIA loops $c^{(111)}_n$}
\label{sec:ME_SIA}
%-------------------------------------------------------------------
The glissile population obeys the full SIA master equation, augmented by the
$\tfrac12\langle111\rangle\!\to\!\langle100\rangle$ coupling (P9). Cascades feed
this population alone, $G_n\equiv G_n^{(111)}$. The junction reaction of two
comparable $\tfrac12\langle111\rangle$ loops removes both reactants from this
population (ordinary coalescence loss) but \emph{redirects} a size-comparability
fraction $\varphi_{\rm junc}$ of the coalescence \emph{product} into the
$\langle100\rangle$ population; the surviving fraction $(1-\varphi_{\rm junc})$
stays glissile. Hence the coalescence gain below carries the
$(1-\varphi_{\rm junc})$ weight, with the complementary product appearing in
Eq.~\eqref{eq:ME_SIA100}.
\begin{equation}\label{eq:ME_SIA}
\begin{split}
\frac{dc^{(111)}_n}{dt} = {} &
  \underbrace{G_n}_{\text{production}}
  + \underbrace{
     \frac{1}{2}\!\!\sum_{n'=1}^{\min(n-1,\,i_{\rm mobile})}\!\!
     \bigl[1-\varphi_{\rm junc}(n',n{-}n')\bigr]\,
     \mathcal{K}_{n',n-n'}^{ii}\,c^{(111)}_{n'}\,c^{(111)}_{n-n'}
     }_{\text{SIA--SIA coalescence gain, glissile product}}
  \\[2pt]
  &- \underbrace{
     c^{(111)}_n\!\sum_{n'=1}^{i_{\rm mobile}}
     \mathcal{K}_{n,n'}^{ii}\,c^{(111)}_{n'}
     }_{\text{SIA--SIA coalescence loss } (I_n+I_{n'}\to I_{n+n'})}
  + \underbrace{
     \varepsilon_i(n{+}1)\,c^{(111)}_{n+1}
   - \varepsilon_i(n)\,c^{(111)}_n\,\delta_{n \geq 2}
     }_{\text{thermal emission (SIA from loops, P5)}}
  \\[2pt]
  &+ \underbrace{
     \sum_{m'=1}^{v_{\rm mobile}}
     \mathcal{K}_{m',n+m'}^{vi}\,c_{m',0}\,c^{(111)}_{n+m'}
     }_{\text{V--I annihilation gain } (V_{m'}+I_{n+m'}\to I_n)}
  - \underbrace{
     \sum_{m'=1}^{v_{\rm mobile}}
     \mathcal{K}_{m',n}^{vi}\,c_{m',0}\,c^{(111)}_n
     }_{\text{V--I annihilation loss } (V_{m'}+I_n\to I_{n-m'})}
  \\[2pt]
  &- \underbrace{
     \sum_{m=1}^{N_V}
     \mathcal{K}_{n,m}^{(\mathrm{cav})}\,c^{(111)}_n\,\bar{c}_m
     }_{\text{SIA cluster $\to$ cavity absorption (P2/P6)}}
  - \underbrace{
     \mathcal{D}_i(n)\,c^{(111)}_n
     }_{\text{fixed-sink loss (mobile clusters only, P4)}}
  \\[2pt]
  &+ \underbrace{
     \delta_{n=1}\sum_{m=1}^{N_V}\sum_{\ell=1}^{\ell_{\max}(m)}
     \Gamma_{\mathrm{TM}}(m,\ell)\,c_{m,\ell}
     }_{\text{trap-mutation SIA source at }n{=}1\text{ (P7)}}
  \;+\; \left.\frac{dc^{(111)}_n}{dt}\right|_{\mathrm{conv}}^{\!\!\!(\mathrm{P9})},
\end{split}
\end{equation}
where the cavity rate constant $\mathcal{K}^{(\mathrm{cav})}_{n,m}$ is
\begin{equation}\label{eq:K_cav}
    \mathcal{K}_{n,m}^{(\mathrm{cav})} =
    \begin{cases}
        A_{\mathrm{sph}}\,m^{1/3}\,\omega_i^{\mathrm{eff}},
        & 1 \leq n \leq 3\;\text{(3D, P2),} \\[5pt]
        \displaystyle
        \frac{A_{\mathrm{sph}}\,m^{1/3}\,\omega_n^{1D}}
             {1 + B_{\mathrm{rot}}\,\hat{L}^2\,m^{-1/3}},
        & 4 \leq n \leq i_{\rm mobile}\;\text{(1D/3D, P6),} \\[5pt]
        0, & n > i_{\rm mobile}\;\text{(sessile)},
    \end{cases}
\end{equation}
and the glissile-side conversion coupling collects the unary depletion, the
$\langle100\rangle$ absorption capture, and the monomer exchange with the
sessile ladders (growth consumes a $\tfrac12\langle111\rangle$ monomer; shrink
of sub-onset loops and thermal emission return one):
\begin{equation}\label{eq:conv_111}
\begin{split}
\left.\frac{dc^{(111)}_n}{dt}\right|_{\mathrm{conv}}^{\!\!\!(\mathrm{P9})}
   = {} & -\,\underbrace{\Gamma_{\rm uni}(n)\,c^{(111)}_n}_{\text{unary reorientation}}
   \;-\; \underbrace{c^{(111)}_n\!\sum_{m\geq n_{\rm conv}} K^{\rm abs}_{m,n}\,c^{(100)}_m}_{\text{absorbed by sessile }\langle100\rangle}
   \\[2pt]
   & +\,\delta_{n=1}\!\sum_{n'\geq n_{\rm conv}}\!\Bigl[
       \varepsilon^{100}(n')\,c^{(100)}_{n'}
     - \mathcal{K}^{100}_{g}(n')\,c^{(100)}_{n'}\,c^{(111)}_{1}\Bigr]
   \;+\; \underbrace{\mathcal{K}^{100}_{s}(n_{\rm conv})\,c^{(100)}_{n_{\rm conv}}\,c_{1,0}\,\delta_{n=n_{\rm conv}-1}}_{\text{shrink dissolves to }\tfrac12\langle111\rangle}.
\end{split}
\end{equation}

%-------------------------------------------------------------------
\paragraph{Defect Class 1b: sessile $\langle100\rangle$ SIA loops $c^{(100)}_n$}
\label{sec:ME_SIA100}
%-------------------------------------------------------------------
The $\langle100\rangle$ population has no cascade source and no self-coalescence
(its loops are immobile, $D^{100}_n\equiv0$). It is created by the three
conversion channels (P9) and subsequently grows or shrinks only by exchanging
single point defects with the matrix --- absorbing $\tfrac12\langle111\rangle$
SIA monomers ($\mathcal{K}^{100}_g$), absorbing vacancy monomers
($\mathcal{K}^{100}_s$, an SIA--vacancy annihilation that shrinks the loop), and
emitting SIA monomers back into the $\tfrac12\langle111\rangle$ pool
($\varepsilon^{100}$):
\begin{equation}\label{eq:ME_SIA100}
\begin{split}
\frac{dc^{(100)}_n}{dt} = {} &
  \underbrace{\Gamma_{\rm uni}(n)\,c^{(111)}_n}_{\text{unary }\tfrac12\langle111\rangle_n\to\langle100\rangle_n\ (\text{A})}
  + \underbrace{
     \frac{1}{2}\!\!\sum_{n'+n''=n}\!\!
     \varphi_{\rm junc}(n',n'')\,\mathcal{K}_{n',n''}^{ii}\,
     c^{(111)}_{n'}\,c^{(111)}_{n''}
     }_{\text{junction nucleation }(\text{B})}
  \\[2pt]
  &+ \underbrace{
     \sum_{n'=1}^{i_{\rm mobile}}
     K^{\rm abs}_{n-n',\,n'}\,c^{(100)}_{n-n'}\,c^{(111)}_{n'}
     }_{\text{absorption growth gain }(\langle100\rangle_{n-n'}+I^{(111)}_{n'}\to\langle100\rangle_n)}
  - \underbrace{
     c^{(100)}_n\!\sum_{n'=1}^{i_{\rm mobile}}
     K^{\rm abs}_{n,\,n'}\,c^{(111)}_{n'}
     }_{\text{absorption growth loss}}
  \\[2pt]
  &+ \underbrace{
     \mathcal{K}^{100}_{g}(n{-}1)\,c^{(100)}_{n-1}\,c^{(111)}_1
   - \mathcal{K}^{100}_{g}(n)\,c^{(100)}_n\,c^{(111)}_1
     }_{\text{SIA-monomer capture growth }(\langle100\rangle_n+I^{(111)}_1\to\langle100\rangle_{n+1})}
  \\[2pt]
  &+ \underbrace{
     \mathcal{K}^{100}_{s}(n{+}1)\,c^{(100)}_{n+1}\,c_{1,0}
   - \mathcal{K}^{100}_{s}(n)\,c^{(100)}_n\,c_{1,0}
     }_{\text{vacancy-monomer capture shrinkage }(\langle100\rangle_n+V_1\to\langle100\rangle_{n-1})}
  \\[2pt]
  &+ \underbrace{
     \varepsilon^{100}(n{+}1)\,c^{(100)}_{n+1}
   - \varepsilon^{100}(n)\,c^{(100)}_n
     }_{\text{thermal SIA emission (P5)}},
   \qquad n \geq n_{\rm conv}.
\end{split}
\end{equation}
The boundary fluxes are handled so that SIA content is conserved across the
onset: an SIA monomer emitted or a vacancy captured at $n=n_{\rm conv}$ returns
the loop to the $\tfrac12\langle111\rangle$ population
[the $\delta$-terms in Eq.~\eqref{eq:conv_111}], and the corresponding vacancy
monomer consumed by shrinkage is debited from $\dot c_{1,0}$ in
Eq.~\eqref{eq:ME_vac}. Because $\langle100\rangle$ loops are immobile they
contribute no $\mathcal{D}_i$ fixed-sink loss and act \emph{as} biased sinks for
the mobile $\tfrac12\langle111\rangle$ flux.

%-------------------------------------------------------------------
\paragraph{$\tfrac12\langle111\rangle\!\to\!\langle100\rangle$ loop conversion (P9)}
\label{sec:loop-conversion}
Irradiated bcc Fe and F/M steels develop prismatic interstitial loops of two
characters: glissile $\tfrac12\langle111\rangle$ loops, mobile to large sizes,
and sessile square $\langle100\rangle$ loops
\citep{Masters1965,EyreBartlett1965}. Elasticity favours the smaller
$\tfrac12\langle111\rangle$ vector and cascades inject almost exclusively
$\tfrac12\langle111\rangle$ clusters \citep{Marian2002}, yet $\langle100\rangle$
loops dominate at elevated temperature and must therefore be \emph{generated}
in situ. We model this with two additive channels feeding
Eqs.~\eqref{eq:ME_SIA}--\eqref{eq:ME_SIA100}: a thermodynamic single-loop
transformation (Channel~A) carrying the temperature trend \citep{Dudarev2008},
and a kinetic junction route (Channel~B) carrying the dose- and
size-dependence \citep{Marian2002,Arakawa2006}.

\textbf{Loop free energies.}
\label{sec:loop-energies}
The free energy of a loop of $n$ SIAs and character $X\in\{111,100\}$ follows
the anisotropic-elasticity form of \cite{Dudarev2008},
\begin{equation}
  E_l^{X}(n,T) = P_X(n)\!\left[\,
    \hat F_X(T)\,\ln\!\frac{4R^\ast_X(n)}{e\,\delta}
    + F_\delta^{X}(T) + F_c^{X}\,\right],
  \label{eq:loop-energy}
\end{equation}
with $P_X$ the perimeter, $R^\ast_X$ the equivalent-circle radius, $\delta$ the
core cutoff, $\hat F_X$ the prelogarithmic energy factor, and core constants in
Table~\ref{tab:core-constants}. The geometry follows from the platelet area
$A=n\Omega/b_X$,
\begin{equation}
  R^\ast_X(n)=\sqrt{A/\pi},\quad
  P_{111}(n)=6\sqrt{\tfrac{2A}{3\sqrt3}}\ \text{(hexagon)},\quad
  P_{100}(n)=4\sqrt{A}\ \text{(square)},
  \label{eq:loop-geometry}
\end{equation}
and the pure-edge $\langle100\rangle[100]$ prelogarithmic factor is analytic and
\emph{vanishes} as the shear constant $c'=\tfrac12(c_{11}-c_{12})$ softens:
\begin{equation}
  \hat F_{001}([100]) = \frac{a^2}{4\pi}(c_{11}{+}c_{12})
  \!\left[\frac{c_{44}(c_{11}{-}c_{12})}{c_{11}(c_{11}{+}c_{12}{+}2c_{44})}\right]^{1/2}
  \;\xrightarrow{\;c'\to0\;}\;0.
  \label{eq:F100}
\end{equation}
The softening is driven by spin fluctuations approaching the $\alpha$--$\gamma$
transition,
\begin{equation}
  c'(T)\simeq 56.8\,(1-T/T_c)^{1/2}\ \text{GPa},\qquad T_c=1185\ \text{K},
  \label{eq:cprime}
\end{equation}
so the $\langle100\rangle$ elastic self-energy collapses at high $T$ while the
$\tfrac12\langle111\rangle$ energy does not.

\textbf{Channel A: single-loop thermodynamic transformation.}
\label{sec:channel-A}
A loop of size $n$ reorients once $\langle100\rangle$ is lower in energy, with
per-loop driving force
\begin{equation}
  \Delta F(n,T) = E_l^{111}(n,T) - E_l^{100}(n,T) = P(n)\,\Delta f(T),
  \label{eq:driving-force}
\end{equation}
parametrised through a per-length difference in which only the
$\langle100\rangle$ prelogarithmic part softens, with exponent fixed at $1/4$ by
$\hat F_{001}\propto(1-T/T_c)^{1/4}$ from Eq.~\eqref{eq:F100}:
\begin{equation}
  \boxed{\;f_{100}(T)=f_{100}^{\rm core}+f_{100}^{\rm pre}\,(1-T/T_c)^{1/4},
  \qquad \Delta f(T) = f_{111} - f_{100}(T)\;}
  \label{eq:delta-f}
\end{equation}
The single fit quantity is the crossover temperature $T^\ast$ (where
$\Delta f=0$), chosen so conversion turns on near $350\,^\circ$C and is strongly
favoured by $550\,^\circ$C, reproducing the three stability regions of
\cite{Dudarev2008}; conversion is admissible for $T>T^\ast$, and the onset size
$n_{\rm conv}$ is the smallest $n$ with $\Delta F(n,T)>0$. The transformation is
thermally activated and one-way; propagating segment-by-segment
\citep{Marian2002}, its barrier scales with the perimeter,
\begin{equation}
  \boxed{\;\Gamma_{\rm uni}(n,T) = \nu_0\,
    \exp\!\left[-\frac{E_a^{\rm uni}(n)}{k_BT}\right]
    \max\!\left[0,\;1-\exp\!\left(-\frac{\Delta F(n,T)}{k_BT}\right)\right],\qquad
    E_a^{\rm uni}(n) = E_a^{0} + \gamma_a\,\frac{P(n)}{b_{111}}\;}
  \label{eq:gamma-uni}
\end{equation}
so the gating factor favours large loops while the size-dependent barrier
suppresses them, giving a preferred conversion-size window that reproduces the
spontaneous transformation of \emph{small} loops on heating \citep{Arakawa2006}
and leaves large loops to convert via Channel~B.

\textbf{Channel B: junction nucleation and absorption growth.}
\label{sec:channel-B}
Two comparable mobile $\tfrac12\langle111\rangle$ loops that collide can react
through the junction \citep{EyreBullough1965,Marian2002}
\begin{equation}
  \tfrac12[111] + \tfrac12[1\bar1\bar1] \rightarrow [100],
  \label{eq:junction}
\end{equation}
forming a single $\langle100\rangle$ loop; this occurs only for similar-size
partners (otherwise the smaller simply rotates into the larger, i.e.\ ordinary
coalescence). We branch the same-polarity kernel $\mathcal{K}^{ii}$
(Eq.~\ref{eq:Kii}) with a size-comparability fraction
\begin{equation}
  \boxed{\;\varphi_{\rm junc}(n,n') = \varphi_{\max}\,P_{\rm s}(T)\,
  \exp\!\left[-\frac{\big(\ln(n/n')\big)^2}{2\sigma_s^2}\right]
  \Theta\!\big(\min(n,n')\ge n_{j,\min}\big)\;}
  \label{eq:phi-junc}
\end{equation}
so that the product fraction $\varphi_{\rm junc}$ enters the $\langle100\rangle$
population [Eq.~\eqref{eq:ME_SIA100}] and the complement $(1-\varphi_{\rm junc})$
remains glissile [Eq.~\eqref{eq:ME_SIA}], with $P_{\rm s}(T)$ the Marian
two-step success probability and
\begin{equation}
  \mathcal K^{ii}_{n,n'} = \frac{8\pi}{\Omega^{2/3}}
  \big(\xi_n+\xi_{n'}\big)\big(D^{111}_n+D^{111}_{n'}\big),\qquad
  \xi_n=\Big(\tfrac{3n}{8\pi}\Big)^{1/3}.
  \label{eq:Kii}
\end{equation}
Once formed, $\langle100\rangle$ loops are stationary and act as biased sinks
for mobile $\tfrac12\langle111\rangle$ clusters \citep{Marian2002},
\begin{equation}
  [100]_m + \tfrac12\langle111\rangle_n \rightarrow [100]_{m+n},
  \label{eq:absorption}
\end{equation}
the dominant route to TEM-visible $\langle100\rangle$ loops, with cross-character
capture kernel driven by the mobile $\tfrac12\langle111\rangle$ partner
($D^{100}_m\approx0$),
\begin{equation}
  \boxed{\;K^{\rm abs}_{m,n} = \frac{8\pi}{\Omega^{2/3}}
  \big(\xi_m+\xi_n\big)\big(D^{100}_m+D^{111}_n\big)
  \approx \frac{8\pi}{\Omega^{2/3}}\big(\xi_m+\xi_n\big)D^{111}_n\;}.
  \label{eq:Kabs}
\end{equation}
The sessile point-defect kernels of Eq.~\eqref{eq:ME_SIA100} are the
loop-absorption forms specialised to a $\langle100\rangle$ target:
$\mathcal{K}^{100}_g(n)=A_{\rm loop}n^{1/2}Z_i^{\rm loop}\omega_i^{\rm eff}$
(SIA-monomer capture), $\mathcal{K}^{100}_s(n)=A_{\rm loop}n^{1/2}\omega_v^{\rm eff}$
(vacancy-monomer capture), and $\varepsilon^{100}(n)$ the detailed-balance
emission rate against $\mathcal{K}^{100}_g$ with $\langle100\rangle$ binding
energy $E_b^{100}(n)=A_{100}\,n^{-B_{100}}$.

\textbf{Conservation of SIA content.}
Every term in Eqs.~\eqref{eq:conv_111}--\eqref{eq:ME_SIA100} conserves the total
SIA inventory $\sum_n n\,(c^{(111)}_n + c^{(100)}_n)$: the unary and junction
channels relabel content at fixed (or summed) size, the absorption channel
carries the summed size into the product, and the growth/emission ladders only
exchange single SIA monomers between the two populations. The only net SIA loss
from the $\langle100\rangle$ population is the vacancy-capture shrinkage, which
is a Frenkel-pair annihilation accounted for symmetrically in $\dot c_{1,0}$.
The Frenkel-pair diagnostic of Section~\ref{sec:FP_balance} is therefore
unaffected by the split, provided $S_I$ sums over both characters
[Eq.~\eqref{eq:SI_S}].

\textbf{Input data and parameters.}
\label{sec:loop-conversion-data}
\begin{table}[ht]
\centering
\caption{Zero-temperature core constants for the loop free
energy~\eqref{eq:loop-energy}, after \cite{Dudarev2008} (eV/\AA).}
\label{tab:core-constants}
\begin{tabular}{lcccc}
\toprule
term & $111[11\bar2]$ & $111[\bar110]$ & $001[100]$ & $001[110]$ \\
\midrule
$F_c$ (nonlinear core)   & 0.46 & 0.47 & 0.33 & --- \\
$F_\delta$ (core-traction) & 0.345 & 0.349 & 0.387 & 0.390 \\
\bottomrule
\end{tabular}
\end{table}

\begin{table}[ht]
\centering
\caption{Conversion-model parameters. Values marked ``cal.'' are calibrated
against the measured $\tfrac12\langle111\rangle$ fraction $f_{111}(T,\text{dose})$.}
\label{tab:conversion-params}
\begin{tabular}{llll}
\toprule
symbol & meaning & first-cut value & source \\
\midrule
$T_c$            & $\alpha$--$\gamma$ softening temperature & 1185 K & \cite{Dudarev2008} \\
$c'(T)$          & shear softening law & $56.8(1{-}T/T_c)^{1/2}$ GPa & \cite{Dudarev2008} \\
$\delta$         & dislocation core cutoff & $0.4$ nm & \cite{Dudarev2008} \\
$f_{111},f_{100}^{0}$ & per-length loop energies & $0.3$--$0.5$ eV/\AA & Tab.~\ref{tab:core-constants} \\
softening exp.   & exponent in $\Delta f(T)$ & $1/4$ (fixed) & \cite{Dudarev2008}, Eq.~\eqref{eq:F100} \\
$T^\ast$         & conversion crossover temperature & cal.\ ($\in[350,550]\,^\circ$C) & \cite{Dudarev2008} \\
$\nu_0$          & attempt frequency & $10^{13}\,\mathrm{s}^{-1}$ & Debye \\
$E_a^{0}$        & unary barrier offset & $0.5$--$2.0$ eV & \cite{Marian2002} \\
$\gamma_a$       & unary barrier size slope & cal. & \cite{Arakawa2006} \\
$\varphi_{\max}$ & peak junction yield ($n{=}n'$) & $0.5$--$1.0$ & \cite{Marian2002} \\
$\sigma_s$       & log-size tolerance & $0.3$--$0.5$ & \cite{Marian2002} \\
$n_{j,\min}$     & minimum junction size & $\mathcal O(10)$ SIAs & \cite{Marian2002} \\
$A_{100},B_{100}$& $\langle100\rangle$ binding-energy fit & $0.716,\ 0.358$ & Table~18 (DFT) \\
$n_{\rm conv}$   & $\langle100\rangle$ onset size & \emph{derived} ($\Delta F{>}0$) & Eq.~\eqref{eq:driving-force} \\
\bottomrule
\end{tabular}
\end{table}

%-------------------------------------------------------------------
\paragraph{Defect Class 2: Vacancy clusters $c_{m,\ell}$}
\label{sec:ME_vac}
%-------------------------------------------------------------------
Voids ($\ell=0$) and bubbles ($\ell\geq 1$) share a unified $(m,\ell)$ grid.
Spherical cavities are unbiased ($Z_i^{\mathrm{cav}}=Z_v^{\mathrm{cav}}=1$). The
v--v coalescence terms are active only for $m\leq v_{\rm mobile}$; SIA-induced
shrinkage spans all mobile $\tfrac12\langle111\rangle$ SIA sizes through
$\mathcal{K}^{(\mathrm{cav})}_{n,m}$. The vacancy monomer additionally feeds the
sessile $\langle100\rangle$ shrinkage ladder of Eq.~\eqref{eq:ME_SIA100}; that
loss is the last term below.
\begin{equation}\label{eq:ME_vac}
\begin{split}
\frac{dc_{m,\ell}}{dt} = {} &
  \underbrace{G_{m,\ell}}_{\text{production}}
  + \underbrace{
     \frac{1}{2}\sum_{m'=1}^{\min(m-1,\,v_{\rm mobile})}
     \mathcal{K}_{m',m-m'}^{vv}\,c_{m',0}\,c_{m-m',\ell}\,
     \delta_{m \leq 2\,v_{\rm mobile}}
     }_{\text{V--V coalescence gain } (V_{m'}+V_{m-m'}\to V_m)}
  \\[2pt]
  &- \underbrace{
     c_{m,\ell}\sum_{m'=1}^{v_{\rm mobile}}
     \mathcal{K}_{m,m'}^{vv}\,c_{m',0}
     }_{\text{V--V coalescence loss } (V_m+V_{m'}\to V_{m+m'})}
  + \underbrace{
     \varepsilon_v(m{+}1,\ell)\,c_{m+1,\ell}
   - \varepsilon_v(m,\ell)\,c_{m,\ell}
     }_{\text{thermal emission (vacancy from cavities, P5)}}
  \\[2pt]
  &+ \underbrace{
     \sum_{n=1}^{\min(i_{\rm mobile},\,N_V-m)}
     \mathcal{K}_{n,m+n}^{(\mathrm{cav})}\,c^{(111)}_n\,c_{m+n,\ell}
     }_{\text{SIA-induced shrinkage gain } (I_n + V_{m+n}\to V_m)}
  - \underbrace{
     \sum_{n=1}^{i_{\rm mobile}}
     \mathcal{K}_{n,m}^{(\mathrm{cav})}\,c^{(111)}_n\,c_{m,\ell}
     }_{\text{SIA-induced shrinkage loss } (I_n + V_m\to V_{m-n})}
  \\[2pt]
  &+ \underbrace{
     A_{\mathrm{sph}}\,m^{1/3}\,\omega_h^{\mathrm{eff}}\,
     c_h\,c_{m,\ell-1}\,\delta_{\ell \geq 1}
   - A_{\mathrm{sph}}\,m^{1/3}\,\omega_h^{\mathrm{eff}}\,
     c_h\,c_{m,\ell}\,\delta_{\ell < \ell_{\max}}
     }_{\text{He absorption by cavities (P2, He-side)}}
  \\[2pt]
  &+ \underbrace{
     \varepsilon_h(m,\ell{+}1)\,c_{m,\ell+1}
   - \varepsilon_h(m,\ell)\,c_{m,\ell}\,\delta_{\ell \geq 1}
     }_{\text{He emission from cavities (P5, He-side)}}
  + \underbrace{
     \Gamma_{\mathrm{TM}}(m{-}1,\ell)\,c_{m-1,\ell}
   - \Gamma_{\mathrm{TM}}(m,\ell)\,c_{m,\ell}
     }_{\text{trap mutation (P7)}}
  \\[2pt]
  &+ \underbrace{
     b_0\,\dot{\phi}\,(\ell{+}1)\,c_{m,\ell+1}
   - b_0\,\dot{\phi}\,\ell\,c_{m,\ell}
     }_{\text{radiation re-solution (P8)}}
  - \underbrace{
     \mathcal{D}_v\,c_{m,\ell}\,\delta_{m \leq v_{\rm mobile}}
     }_{\text{fixed-sink loss (mobile cavities only, P4)}}
  \\[2pt]
  &- \underbrace{
     \delta_{m=1}\,\delta_{\ell=0}\,c_{1,0}\!
     \sum_{n\geq n_{\rm conv}} \mathcal{K}^{100}_{s}(n)\,c^{(100)}_n
     }_{\text{vacancy monomer consumed by }\langle100\rangle\text{ shrinkage (P9)}}.
\end{split}
\end{equation}

%-------------------------------------------------------------------
\paragraph{Free helium $c_h$}
\label{sec:ME_He}
%-------------------------------------------------------------------
He--He coalescence is excluded, and helium interacts with the vacancy-cluster
grid only through mono-He trapping and emission; the $\langle100\rangle$
population carries no helium and does not appear here.
\begin{equation}\label{eq:ME_He}
\begin{split}
\frac{dc_h}{dt} = {} &
  \underbrace{G_{\mathrm{He}}}_{\text{production}}
  - \underbrace{
     \sum_{m=1}^{N_V}
     A_{\mathrm{sph}}\,m^{1/3}\,\omega_h^{\mathrm{eff}}\,c_h
     \sum_{\ell=0}^{\ell_{\max}(m)-1} c_{m,\ell}
     }_{\text{He absorption by cavities (P2)}}
  \\[2pt]
  &- \underbrace{
     \mathcal{D}_h\,c_h
     }_{\text{fixed-sink loss (P4)}}
  + \underbrace{
     \sum_{m=1}^{N_V}\sum_{\ell=1}^{\ell_{\max}(m)}
     \varepsilon_h(m,\ell)\,c_{m,\ell}
     }_{\text{He emission (P5)}}
  + \underbrace{
     \sum_{m=1}^{N_V}\sum_{\ell=1}^{\ell_{\max}(m)}
     b_0\,\dot{\phi}\,\ell\,c_{m,\ell}
     }_{\text{radiation re-solution (P8)}}.
\end{split}
\end{equation}

%-------------------------------------------------------------------
\paragraph{Exact conservation identities}
%-------------------------------------------------------------------
The master equations~\eqref{eq:ME_SIA}--\eqref{eq:ME_He} satisfy two exact
conservation identities. Their derivation rests on three facts built into the
\textsc{RadCluster} code:
\begin{enumerate}
  \item[(i)] cascade production deposits SIA and vacancy atoms in a balanced
        Frenkel-pair ratio (and feeds only the $\tfrac12\langle111\rangle$
        character),
        $\sum_{n=1}^{N_I} n\,P_n^{(i)}
         = \sum_{m=1}^{N_V} m\,P_m^{(v)} = \eta G$;
  \item[(ii)] every mutual reaction destroys $\min(n,m)$ SIA atoms \emph{and}
        $\min(n,m)$ vacancy atoms with the residual fragment surviving in its
        population; this holds for the mobile $\tfrac12\langle111\rangle$--vacancy
        annihilations \emph{and} for the sessile $\langle100\rangle$
        vacancy-capture shrinkage (each event removes one SIA from
        $c^{(100)}$ and one vacancy from $c_{1,0}$), so the two destruction
        sums cancel identically in $d(S_I-S)/dt$; the
        $\tfrac12\langle111\rangle\!\leftrightarrow\!\langle100\rangle$
        conversion edges only relabel SIA content and drop out;
  \item[(iii)] only mobile clusters with $n \le i_{\rm mobile}$ (glissile SIAs)
        or $m \le v_{\rm mobile}$ (vacancies) reach fixed sinks. The sessile
        $\langle100\rangle$ loops are sinks, not sink-bound species, and carry no
        $\mathcal{D}_i$ term; the rate-equation losses are summed over the full
        mobile range, not the monomer alone.
\end{enumerate}

\paragraph{Cumulative loss to fixed sinks}
The total fixed-sink rate for species $\alpha\in\{i,v,h\}$ collects all three
sink channels,
$\mathcal{D}_\alpha(n) = \mathcal{D}_\alpha^d(n) +
 \mathcal{D}_\alpha^{gb}(n) + \mathcal{D}_\alpha^p(n)$, and vanishes outside the
mobile range. Define the cumulative atomic content delivered to fixed sinks
(only the glissile SIAs are mobile):
\begin{align}
    \Delta J_i^{\mathrm{fix}}(t)
    &\equiv \int_0^t \sum_{n=1}^{i_{\rm mobile}}
        n\,\mathcal{D}_i(n)\,c^{(111)}_n(t')\,dt' \;\ge\; 0,
    \label{eq:DeltaJi_fix}\\[3pt]
    \Delta J_v^{\mathrm{fix}}(t)
    &\equiv \int_0^t \sum_{m=1}^{v_{\rm mobile}}
        m\,\mathcal{D}_v\,\bar{c}_m(t')\,dt' \;\ge\; 0,
    \label{eq:DeltaJv_fix}\\[3pt]
    \Delta J_h^{\mathrm{fix}}(t)
    &\equiv \int_0^t\!\!\left[
        \mathcal{D}_h\,c_h
        + \mathcal{D}_v
          \sum_{m=1}^{v_{\rm mobile}}\sum_{\ell}\ell\,c_{m,\ell}
      \right]\!dt' \;\ge\; 0.
    \label{eq:DeltaJh_fix}
\end{align}
These three integrals are appended to the state vector and evolved as auxiliary
ODEs, so that they share the integrator's local error control with the physical
concentrations.

\paragraph{Frenkel-pair balance}
\label{sec:FP_balance}
With total inventories summed over \emph{both} SIA characters,
\begin{equation}\label{eq:SI_S}
    S_I(t) \equiv \sum_{n=1}^{N_I} n\,\bigl[c^{(111)}_n(t)+c^{(100)}_n(t)\bigr],
    \qquad
    S(t)   \equiv \sum_{m=1}^{N_V} m\,\bar{c}_m(t),
\end{equation}
facts (i)--(iii) combine into the exact Frenkel-pair conservation law
\begin{equation}\label{eq:mass_FP}
  S(t) - S_I(t)
  \;-\; \Delta J_i^{\mathrm{fix}}(t)
  \;+\; \Delta J_v^{\mathrm{fix}}(t)
  \;=\; S(0) - S_I(0)
  \;=\; \mathrm{const.}
\end{equation}
For a defect-free initial condition, this is the \emph{swelling identity}
\begin{equation}\label{eq:swelling_identity}
    S(t) \;=\; S_I(t)
          \;+\; \Delta J_i^{\mathrm{fix}}(t)
          \;-\; \Delta J_v^{\mathrm{fix}}(t),
\end{equation}
i.e.\ the total vacancy content of all cavities equals the SIA inventory in
loops (both characters) plus the \emph{net} atomic flux delivered to fixed
sinks. Differentiating gives the rate form
\begin{equation}\label{eq:dSdt}
    \frac{dS}{dt}
    \;=\; \frac{dS_I}{dt}
    \;+\; \sum_{n=1}^{i_{\rm mobile}} n\,\mathcal{D}_i(n)\,c^{(111)}_n
    \;-\; \sum_{m=1}^{v_{\rm mobile}} m\,\mathcal{D}_v\,\bar{c}_m.
\end{equation}
Once loops saturate ($dS_I/dt \approx 0$), swelling grows at the rate at which
the biased dislocation channel ($Z_i^d \neq Z_v^d$) drains the surplus SIA
inventory faster than the vacancy inventory; the unbiased grain-boundary and
precipitate channels contribute equally to both sums and cancel in
$dS/dt - dS_I/dt$ only insofar as the mobile SIA and vacancy populations are
equal.

\paragraph{Helium conservation}
The total helium inventory is the sum of free and trapped helium,
\begin{equation}
    S_{\mathrm{He}}(t)
    \;\equiv\; c_h(t)
    \;+\; \sum_{m=1}^{N_V}\sum_{\ell=1}^{\ell_{\max}(m)}
              \ell\,c_{m,\ell}(t).
\end{equation}
Trap mutation (P7), radiation re-solution (P8), and thermal He emission (P5h)
only redistribute helium between $c_h$ and the trapped pool, so their
contributions to $dS_{\mathrm{He}}/dt$ cancel pairwise. The only net helium
sinks are free helium captured at fixed sinks ($\mathcal{D}_h\,c_h$) and helium
carried to fixed sinks by mobile cavities ($m \le v_{\rm mobile}$); immobile
bubbles retain their helium. Hence
\begin{equation}\label{eq:mass_He}
    \frac{dS_{\mathrm{He}}}{dt}
    \;=\; G_{\mathrm{He}}
    \;-\; \mathcal{D}_h\,c_h
    \;-\; \mathcal{D}_v
          \sum_{m=1}^{v_{\rm mobile}}\sum_{\ell} \ell\,c_{m,\ell},
\end{equation}
i.e.\ $dS_{\mathrm{He}}/dt = G_{\mathrm{He}} - d\Delta J_h^{\mathrm{fix}}/dt$.
In particular, $S_{\mathrm{He}}$ is \emph{not} conserved against
$\int G_{\mathrm{He}}\,dt'$ alone --- the fixed-sink loss must be tracked
explicitly. Both loss terms are typically small (mobile clusters carry little
He) but they are non-zero and must appear in the diagnostic below for it to
converge to zero.

\paragraph{Numerical accuracy diagnostics}
\label{sec:conservation_check}
The swelling identity~\eqref{eq:swelling_identity} and the helium
balance~\eqref{eq:mass_He} provide two inexpensive a-posteriori checks, neither
of which requires post-processing because the auxiliary
integrals~\eqref{eq:DeltaJi_fix}--\eqref{eq:DeltaJh_fix} are already state
variables:
\begin{equation}\label{eq:delta_FP}
    \delta_{\mathrm{FP}}(t)
    = \frac{\bigl|\,S(t) - S_I(t)
            - \Delta J_i^{\mathrm{fix}}(t)
            + \Delta J_v^{\mathrm{fix}}(t)\,\bigr|}
           {S(t) + S_I(t)
            + \Delta J_i^{\mathrm{fix}}(t)
            + \Delta J_v^{\mathrm{fix}}(t)},
\end{equation}
\begin{equation}\label{eq:delta_He}
    \delta_{\mathrm{He}}(t)
    = \frac{\Bigl|\,
        S_{\mathrm{He}}(t)
        + \Delta J_h^{\mathrm{fix}}(t)
        - \displaystyle\int_0^t G_{\mathrm{He}}\,dt'
        - S_{\mathrm{He}}(0)\,\Bigr|}{
        S_{\mathrm{He}}(t)
        + \Delta J_h^{\mathrm{fix}}(t)
        + \displaystyle\int_0^t G_{\mathrm{He}}\,dt'}.
\end{equation}
Both should remain below $\sim 10^{-6}$ in a correctly implemented code; values
above $10^{-3}$ almost always signal a coding error, missing or incorrect
reactions, or an overly loose integrator tolerance. We have verified that the
two-population split preserves these bounds: with the
$\tfrac12\langle111\rangle\!\to\!\langle100\rangle$ conversion active,
$\delta_{\mathrm{FP}}$ remains at the $10^{-7}$ level in the bin-moment
representation (and at the round-off level $\sim10^{-14}$ in the fully discrete
representation), confirming that the conversion edges and the sessile
$\langle100\rangle$ ladders are SIA-content conserving by construction. Despite
the importance of these conservation identities, they have seldom been discussed
or tested in the cluster dynamics literature.
